\documentclass[a4j,11pt,dvipdfmx]{jsarticle}
\usepackage{float}
\usepackage{array}
\usepackage{titlesec}
\usepackage[dvipdfmx]{graphicx}
\usepackage{graphicx}
\usepackage{url}

\titleformat*{\section}{\bfseries}
\titleformat*{\subsection}{\bfseries}

\makeatletter
  \renewcommand{\section}{%
    \@startsection{section}{1}{\z@}%
    {0.4\Cvs}{0.1\Cvs}%
    {\normalfont\headfont\raggedright}}
\makeatother

\makeatletter
  % sectionの下マージンを小さく
  \renewcommand{\subsection}{%
    \@startsection{subsection}{1}{\z@}%
    {0.1\Cvs}{0.1\Cvs}%
    {\normalfont\headfont\raggedright}}
\makeatother

\renewcommand{\thesubsection}{\thesection-\arabic{subsection}.}


%---------------------------------------------------
% ページの設定
%---------------------------------------------------
\setlength{\textwidth}{160truemm}
\setlength{\textheight}{250truemm}
\setlength{\topmargin}{-4.5truemm}
\setlength{\oddsidemargin}{0.5truemm}
\pagestyle{empty}
\setlength{\headheight}{0truemm}
\setlength{\parindent}{1zw}
\newcolumntype{b}{!{\vrule width 1pt}}
\newcommand{\bhline}{\noalign{\hrule height 1pt}}   


\begin{document}
提出日: \today
\begin{center}
\huge{システム評価のためのテスト項目}\\
\large{24G1051 久峩 丈}
\end{center}
\begin{flushright}
\end{flushright}
% \begin{table}[H]
%   \centering
%   \begin{tabular}{bp{15.5truecm}b}
%   \bhline
%   \large{\bf{}}
%   \\ \bhline
% \end{tabular}
% \end{table}


%%%%%%%%%%%%%%%%%%%%%%%%%%%%%%%%%%%%%%%%%%%%%%%%%
%%%%%%%%%%%%%%%%%%%%%%%%%%%%%%%%%%%%%%%%%%%%%%%%%
%%%%%%%%%%%%%%%%%%%%%%%%%%%%%%%%%%%%%%%%%%%%%%%%%
\section{部門における評価指標・テスト実施方法}\label{sec:intro}
\begin{table}[H]
\vspace{-1em}
\centering
\begin{tabular}{bp{15.5truecm}b}
\bhline
\subsection{部門名}
音通信

\\ \bhline
\subsection{評価内容概要}
\begin{itemize}
    \item 通信精度，通信速度，通信距離
\end{itemize}
本システムを通信精度，通信速度，通信距離の3点で評価する．
各評価の目標，方法は以下の通りである．
\\ \bhline
\subsection{評価指標}
\begin{itemize}
    \item 通信精度: $BER = 10^{-1}$
    \item 通信速度: bps $\geq$ 10
    \item 通信距離: 10 m以上
\end{itemize}
各評価指標について，それぞれ以上のものを目標として評価する．
また，BERはビット誤り率である．

\\ \bhline
\subsection{評価実施方法・手順}
通信精度
\begin{enumerate}
    \item 初めに送信する文字を外部ツールを使って4進数に変換しておく
    \item 受信した4進数生データと1．を比較してBERを算出する
\end{enumerate}
通信速度
\begin{enumerate}
    \item 初めに送信する文字を外部ツールを使って2進数に変換しておく
    \item 音が鳴り始めたら時間を測定開始
    \item 受信終了（"----- End receiving -----"のシリアル出力）したら測定終了
    \item 測定時間と1．からbpsを計算する
\end{enumerate}
通信距離
\begin{enumerate}
    \item スピーカーとマイクの距離を0cmから1cmずつ離していく
    \item BER $\geq$ 10を保てる距離で最大のものを通信可能距離とする
\end{enumerate}
\\ \bhline

\end{tabular}
\end{table}

%%%%%%%%%%%%%%%%%%%%%%%%%%%%%%%%%%%%%%%%%%%%%%%%%
%%%%%%%%%%%%%%%%%%%%%%%%%%%%%%%%%%%%%%%%%%%%%%%%%
%%%%%%%%%%%%%%%%%%%%%%%%%%%%%%%%%%%%%%%%%%%%%%%%%

\end{document}
